
\documentclass[12pt,a4paper]{article}

\usepackage[margin=1in]{geometry}
\usepackage{fontspec}
\usepackage{xcolor}
\usepackage{amsmath}
\usepackage{makeidx}
\usepackage{fancyhdr}
\usepackage{setspace}
\usepackage{url}

\setmainfont{Times New Roman}
\setsansfont{Arial}
\setmonofont{Courier New}

\setlength{\emergencystretch}{3em}

\makeindex

\pagestyle{fancy}
\fancyhf{}
\fancyhead[L]{Assignment 6}
\fancyhead[C]{Document Structure}
\fancyhead[R]{LaTeX}
\fancyfoot[C]{\thepage}
\setlength{\headheight}{15pt}

\renewcommand{\thesection}{\Alph{section}}

\begin{document}
	
	\pagecolor{yellow!10}
	\color{black}
	
	\begin{center}
		{\Huge \textbf{Assignment 6}}\\[0.5cm]
		{\Large Document Structure}\\[0.5cm]
	\end{center}
	
	\thispagestyle{empty}
	
	\section{Fonts and Background}
	
	This paragraph demonstrates the use of different fonts and background formatting.
	
	\begin{trivlist}
		\item {\rmfamily This text is written using the serif font Times New Roman.}
		\item {\sffamily This text is written using the sans-serif font Arial.}
		\item {\ttfamily This text is written using the monospaced font Courier New.}
	\end{trivlist}
	
	\index{background}
	
	\section{Spacing}
	
	This is the first line of the paragraph.
	This is the second line of the same paragraph.
	
	\vspace{1cm}
	
	This paragraph begins after a vertical space of one centimetre.
	
	The word antidis-establishmentarianisms demonstrates manual hyphenation.
	
	The word antidisestablishmentarianisms demonstrates automatic hyphenation.
	
	\medskip
	
	LaTeX provides several commands for controlling the spacing between paragraphs and lines. These commands can be used to improve the overall appearance and readability of a document.
	
	\newpage
	
	\setcounter{section}{3}
	
	\section{Numbering}
	
	This section demonstrates custom section numbering using alphabets.
	
	The equation of an ellipse is:
	
	\begin{equation*}
		\frac{x^2}{a^2}+\frac{y^2}{b^2}=1
	\end{equation*}
	
	The equation of a circle is:
	
	\begin{equation}
		x^2+y^2=1
		\label{eq:circle}
	\end{equation}
	
	Equation~\ref{eq:circle} represents a circle with unit radius.
	
	\section{Boxes and Lines}
	
	A savebox can store text and reuse it multiple times.
	
	\newsavebox{\mybox}
	\savebox{\mybox}{This message is written in a savebox.}
	
	\noindent\usebox{\mybox}
	
	\noindent\usebox{\mybox}
	
	\noindent\usebox{\mybox}
	
	\vspace{0.5cm}
	
	The following message is placed inside a paragraph box of width 5 cm:
	
	\noindent
	\parbox{5cm}{
		This message is written inside a paragraph box. The width of this box is five centimetres.
	}
	
	\vspace{0.5cm}
	
	Boxes are useful for highlighting text and controlling the layout of content in a document.
	
	\index{numbering}
	\index{savebox}
	
	\newpage
	
	\section*{Index}
	\addcontentsline{toc}{section}{Index}
	\printindex
	
	\newpage
	
	\section*{References}
	\addcontentsline{toc}{section}{References}
	
	\begin{thebibliography}{9}
		
		\bibitem{notso}
		T. Oetiker,
		\textit{The Not So Short Introduction to LaTeX 2e},
		CTAN, 2021.
		
		\bibitem{overleaf}
		Overleaf,
		\textit{Bibliography Management with BibTeX}.
		Available at:
		\url{https://www.overleaf.com/learn/latex/Bibliography_management_with_bibtex}
		
	\end{thebibliography}
	
\end{document}
